\scanpage{223}
\mspage{114}
$\overline{\underline{\unsure{\mathrm{X})|\mathrm{XI}|\mathrm{XI})}}}$ \uline{Récapitulation}

\deleted{\unsure{4. 2.|4.} Récapitulation}

\uline{1}. soit
% slanted
\margin{Graphes, prégraphes, jauges, multigraphes, multijauges\ldots}
\[
 \Gamma=(\overset{\text{sommets}}S,\overset{\text{cosommets}}{S^\vee},
 \overset{\text{incidence}}R\deleted{\unsure{\ni}}\subset S\times S^\vee) \tag*{(1)}
\]
un ``pré-graphe''. Deux applications
\[
 \varphi:S\longrightarrow S^\vee,\qquad \psi^{\deleted{\unsure{\vee}}}:S^{\deleted{\unsure{\vee}}}\longrightarrow S^\vee
 \tag*{(2)}
\]
sont dites ``associées'' (ou ``$R$-associées'') si on a
\[
 \forall s,t\in S\qquad (s,\psi(t))\in R \quad\Longleftrightarrow\quad
 \deleted{\varphi(s),\psi}(t,\varphi(s))\in R. \tag*{(3)}
\]
\deleted{\unsure{Si}} \deleted{\unsure{i.e.}} \deleted{$\unsure{t\varphi sRt}$}

C'est une relation symétrique en $\varphi,\psi$. Si $R$ est
séparée en $S^\vee$ [i.e.
$\begin{aligned} S^\vee&\longleftrightarrow\mathcal P(S)\\
s^\vee&\longmapsto{}^sR(s^\vee)\end{aligned}$ injective]
alors pour tout $\varphi:S\longrightarrow S^\vee$, il existe au plus une
$\psi:S\longrightarrow S^\vee$ associée -- \deleted{on dit alors que $\varphi$}
\inserted{en tous cas, on} appelle quasi-jauges de $\Gamma$ les $\varphi:S\longrightarrow S^\vee$ telles
qu'il existe une $\psi$ associée.
\deleted{D'où sur l'ens des}
\inserted{Sur l'ens des quasi-jauges, on a donc, dans le cas $\Gamma$ sép. en $S^\vee$, une involution
$(3')\ \varphi\longmapsto{}^s\varphi$.}
\deleted{quasi-jauges, il y a une structure de graphe,
\unsure{définie par} les ``arêtes'' sont les parties de la forme
$\{\varphi,\psi\}$, avec $\varphi$ et $\psi$ associées et sur l'ensemble de
\unsure{celles-ci}. Dans le cas $\Gamma$ séparé en $S^\vee$, une $\psi$ est une}

On appelle \uline{jauge} une $\varphi$ associée à elle $=$,
\unsure{i.e.|i.e. est} telle que
\[
 \forall(s,t)\in S\times S,\qquad (s,\varphi(t))\in R
 \ \Longrightarrow\ \deleted{\unsure{\varphi(s)}}(t,\varphi(s))\in R \tag*{(4)}
\]
ou encore, dans le cas où $\Gamma$ est séparé en $S^\vee$
\[
 \varphi={}^s\varphi .\tag*{(5)}
\]

Supposons que \deleted{l'on ait} $\varphi$ soit une bijection
\inserted{(on dit alors que $\varphi$ est une \uline{jauge stricte})}
\[
 \varphi:S\xrightarrow{\sim}S^\vee\qquad\text{d'où }\varphi'^{\deleted{\unsure{\vee}}}=\varphi^{-1}:S^\vee\xrightarrow{\sim}S
 \tag*{(6)}
\]
et considérons \deleted{l'application} \inserted{l'isom.} de $\underline{(\mathrm{Ens})\times(\mathrm{Ens})}$
\[
 (S,S^\vee)\xrightarrow[(\varphi,\varphi'\overset{\mathrm{déf}}=\deleted{\unsure{\vee}}\varphi^{-1})]{}
 (S'\overset{\mathrm{déf}}=S^\vee,S'{}^\vee\overset{\mathrm{déf}}=S)
 \tag*{(7)}
\]
en sorte que \deleted{\unsure{cette|ce} couple de bijections transforme}
% archive sheet number: 222

